\scanpage{205}
\mspage{79}
Cela précise déjà la définition récurrente pour $\alpha\inserted{=\alpha_0}$ ordinal limite. Notons que le $A_{\alpha_0}$ construit en termes des $A_\beta$ $(\beta<\alpha_0)$ satisfait bien aux conditions que

a) les $\unsure{i_{\alpha_0,\beta}|i_{\beta,\beta}}:A_\beta\unsure{-|\to} A_{\alpha_0}$ sont dans $\widetilde C_{\unsure{\pi|R}}$ \deleted{\unsure{b)}} car $\widetilde C_\pi$ stable par ``limites inductives ordinales strictes'' (cf 1.19. p. 8).

b) $j_{\alpha_0}: A_{\alpha_0}\longrightarrow B$ est dans $C$, car c'est la limite inductive \inserted{filtrante} \deleted{\unsure{directe|d'un}} \unsure{du} système des $A_\beta\to B$ $(\beta<\alpha_0)$ qui sont dans $C$, et $C$ est stable par ces limites.

Il reste donc à préciser comment on passe de $A_\alpha$ à $A_{\alpha+1}$. Cela revient à prouver le

\uline{Lemme 6.4} Soit $f:A\lhook\joinrel\longrightarrow B$ dans $C$, \inserted{telle que $f$ ne soit pas un iso.} Alors il existe un sous-objet $A_1$ de $B$, \uline{distinct de} $A$, tel que $A\unsure{\to|-} A_1$ soit dans $\widetilde{C^\pi}$, et $A_1\to B$ dans $C$.

Si ce lemme est admis, cela signifie que la construction transfinie se continue aussi longtemps que l'on trouve des $A_\alpha$ tels que $A_\alpha\ne B$ i.e. $A_\alpha\hookrightarrow B$ non iso. Mais comme les $A_\alpha$ sont tous différents, si $\pi$ est l'ens des sous-objets de $B$ (on sait que c'est un
% archive sheet number: 204
